\documentclass[a4paper,12pt,twoside]{article}

% Include preamble
% Margins
\usepackage[a4paper, top=2.5cm, bottom=2.5cm, left=2.5cm, right=2.5cm]{geometry}
\setlength{\footskip}{35pt}

% Language and encoding
\usepackage[utf8]{inputenc}
\usepackage[spanish]{babel}

% Font and spacing
\usepackage{setspace}
\onehalfspacing

% Images
\usepackage{graphicx}
\usepackage{float}
\usepackage{svg}

% Tables
\usepackage{array}
\usepackage{booktabs}
\usepackage{colortbl}

% References in APA style
\usepackage{apacite}
\bibliographystyle{apacite}

% Index for elements
\usepackage{tocloft}
\usepackage{listings}
\usepackage{caption}

% Hyperlinks
\usepackage{hyperref}
\hypersetup{
    colorlinks=false,
    linkcolor=blue,
    filecolor=magenta,      
    urlcolor=cyan,
}

% Colors and code
\usepackage{xcolor}
\usepackage{listings}

\lstdefinestyle{darkula}{
    backgroundcolor=\color{black},
    basicstyle=\ttfamily\color{white},
    keywordstyle=\color{orange},
    stringstyle=\color{cyan},
    commentstyle=\color{gray},
    numberstyle=\tiny\color{lightgray},
    numbers=left,
    numbersep=5pt,
    frame=single,
    tabsize=4,
    showstringspaces=false,
    captionpos=b,
    breaklines=true,
    breakatwhitespace=false,
    title=\lstname,
    morekeywords={self}
}

% Rename list of listings and listing captions
\renewcommand{\lstlistlistingname}{Códigos}
\renewcommand{\lstlistingname}{Código}

% Fancy header and footer
\usepackage{fancyhdr}
\setlength{\headheight}{55pt}
\addtolength{\topmargin}{-23pt}
\renewcommand{\headrulewidth}{1pt}
\renewcommand{\footrulewidth}{1pt}

% Title and author
\title{Actividad 1A - 01MBID}
\author{Jahir Gilberth Medina Lopez}
\date{\today}

\def\mytitle{Actividad 1A - 01MBID}
\def\myauthor{Jahir Gilberth Medina Lopez}
\def\mydate{\today}

% Page styles
\fancypagestyle{plain}{
        \fancyhf{}
\fancyhead[L]{
        \includegraphics[width=0.2\textwidth]{assets/images/VIU-Logo.png}
    } % Adjust the path to your header image

        \fancyfoot[LE]{
        \textcolor{white}{\colorbox{orange}{\makebox[1.5em][c]{\thepage}}}
        \hspace{1cm}
        \fontsize{6}{8}\selectfont\MakeUppercase{Máster Universitario en Big Data y Ciencia de Datos | Edición Octubre 2024}
    }
\fancyfoot[RO]{
        \fontsize{9}{11}\selectfont\MakeUppercase{\mytitle}
        \hspace{1cm}
        \textcolor{white}{\colorbox{orange}{\makebox[1.5em][c]{\thepage}}}
    }

    }
\fancypagestyle{toc}{
        \fancyhf{}
\fancyhead[L]{
        \includegraphics[width=0.2\textwidth]{assets/images/VIU-Logo.png}
    } % Adjust the path to your header image

        \fancyfoot[LE]{
        \textcolor{white}{\colorbox{orange}{\makebox[1.5em][c]{\thepage}}}
        \hspace{1cm}
        \fontsize{6}{8}\selectfont\MakeUppercase{Máster Universitario en Big Data y Ciencia de Datos | Edición Octubre 2024}
    }
\fancyfoot[RO]{
        \fontsize{9}{11}\selectfont\MakeUppercase{\mytitle}
        \hspace{1cm}
        \textcolor{white}{\colorbox{orange}{\makebox[1.5em][c]{\thepage}}}
    }

    }
\fancypagestyle{main}{
        \fancyhf{}
\fancyhead[L]{
        \includegraphics[width=0.2\textwidth]{assets/images/VIU-Logo.png}
    } % Adjust the path to your header image

        \fancyfoot[LE]{
        \textcolor{white}{\colorbox{orange}{\makebox[1.5em][c]{\thepage}}}
        \hspace{1cm}
        \fontsize{6}{8}\selectfont\MakeUppercase{Máster Universitario en Big Data y Ciencia de Datos | Edición Octubre 2024}
    }
\fancyfoot[RO]{
        \fontsize{9}{11}\selectfont\MakeUppercase{\mytitle}
        \hspace{1cm}
        \textcolor{white}{\colorbox{orange}{\makebox[1.5em][c]{\thepage}}}
    }

    }

% Ensure text is justified
\usepackage{ragged2e}
\justifying
\raggedbottom


\begin{document}

% Title Page
\thispagestyle{empty}
\begin{center}
    \vspace*{-4.2cm} % Adjust this value if necessary
    \includegraphics[width=0.3\textwidth]{assets/images/VIU-Title.png}\\[1cm] % Adjust the path to your logo
    \vspace*{1cm}
    \textbf{\LARGE Máster Universitario en Big Data y Ciencia de Datos}\\[0.5cm]
    \textit{\MakeUppercase{\mytitle}}\\[2cm]
    \vspace*{3cm}
    \begin{flushleft}
        \textbf{Alumno:} \textit{\MakeUppercase{\myauthor}}\\[0.5cm]
        Edición Octubre 2024 a \textbf{\mydate}
    \end{flushleft}
\end{center}
\pagebreak


\newpage

% Start using the custom page style
\pagestyle{toc}

\tableofcontents
\listoffigures
\listoftables
\lstlistoflistings
\pagebreak

\pagestyle{main}

\section{Introduccion}

\pagebreak

\section{Metodologia}

\begin{table}[H]
    \centering
    \begin{tabular}{|lccccc|}
        \hline
        \rowcolor{orange} \textbf{Operation} & \textbf{Symbol} & \textbf{Example} & & & \\ 
        \hline
        Suma & + & $4 + 4 = 8$ & & & \\
        \hline
        Resta & - & $4 - 2 = 2$ & & & \\
        \hline
        Multiplicación & * & $4 * 4 = 16$ & & & \\
        \hline
    \end{tabular}
    \caption{Random numeric values for different parameters}
    \label{tab:aritmetic_operations}
\end{table}

\pagebreak

\section{Resultados}

\begin{table}[H]
    \centering
    \begin{tabular}{|cccccc|}
        \hline
        \rowcolor{orange} & Temperature & Volume & Size & Age & Total \\
        \hline
        1 & 23.5 & 45.2 & 12.3 & 5 & 86.0 \\
        \hline
        2 & 25.1 & 47.8 & 13.1 & 6 & 92.0 \\
        \hline
        3 & 22.8 & 44.5 & 11.9 & 4 & 83.2 \\
        \hline
        4 & 24.3 & 46.7 & 12.8 & 5 & 89.8 \\
        \hline
        5 & 23.9 & 45.9 & 12.5 & 5 & 87.8 \\
        \hline
    \end{tabular}
    \caption{Random numeric values for different parameters}
    \label{tab:random_values}
\end{table}

\lstinputlisting[language=Python, style=darkula, caption={Example Python Code}]{example.py}

\pagebreak

\section{Discusion}

\pagebreak

\section{Conclusion}

\pagebreak

\bibliography{references}

\end{document}
